\usepackage{hyperref}
\usepackage[utf8]{inputenc}
\usepackage[toc,nopostdot,nogroupskip,nogroupskip,numberedsection=nolabel,automake,acronyms,nonumberlist,nomain]{glossaries-extra}

\makeglossaries

\newacronym{UKF}{UKF}{
    (Unscented Kalman Filter) Сигма-точковий фільтр Калмана
}
\newacronym{EKF}{EKF}{
    (Extended Kalman Filter) Розширений фільтр Калмана
}
\newacronym{MoG2}{MoG2}{
    (Mixture of Gaussians 2) Суміш гауссових коефіцієнтів 2
}
\newacronym{CNN}{CNN}{
    (Convolutional Neural Network) Згорткова нейронна мережа
}
\newacronym{RPN}{RPN}{
    (Region Proposal Network) Мережа пропозиції регіонів
}
\newacronym{CV}{CV}{
    (Computer Vision) Комп'ютерний зір
}
\newacronym{YOLO}{YOLO}{
    (You Only Look Once) Ти дивишся тільки раз
}
\newacronym{IoU}{IoU}{
    (Intersection over Union) Перетин над об'єднанням
}
\newacronym{SORT}{SORT}{
    (Simple Online and Realtime Tracking) Просте онлайн-відстеження та відстеження в режимі реального часу
}
\newacronym{PnP}{PnP}{
    (Perspective-n-Point) Перспектива n-ої точки
}
\newacronym{IMM}{IMM}{
    (Interacting Multiple Model) Взаємодіюча множинна модель
}
\newacronym{Faster R-CNN}{Faster R-CNN}{
    (Faster Region-based Convolutional Neural Network) Швидша згорткова нейронна мережа на основі регіонів
}


%\glsfindwidesttoplevelname
\setglossarystyle{super3col}
\setlength{\glsdescwidth}{.83\textwidth}
%\setabbreviationstyle{short-nolong}

